% ============================================================================
% Fichier  : partie5/P05_C22_cellule.tex
% Titre    : La Cellule de Crise
% Auteur   : MINKA MI NGUIDJOÏ Thierry Emmanuel
% Version  : 1.0 -- novembre 2026
% ============================================================================
% Sources : ISO 22361:2022 ; ISO 22320:2018 (gestion des situations
% d'urgence). Main courante de l'Opération GAROUA : cas fictif.
% ============================================================================

\chapter{La Cellule de Crise}
\label{chap:cellule}
\chaptermark{La cellule de crise}

\epigraph{%
    « En crise, ce qui n'est pas écrit n'a pas été décidé. »%
}{Principe du cours}

\niveauChapitre{Intermédiaire à expert}{Chapitres~\ref{chap:incidents}
et~\ref{chap:crise}.}

\begin{boxobjectifs}
\begin{itemize}
    \item Organiser une cellule de crise à deux niveaux, stratégique et
          opérationnel.
    \item Attribuer les rôles~: direction, coordination, main courante,
          communication, juridique, experts, logistique.
    \item Tenir une main courante exploitable.
    \item Rythmer la crise par des points de situation et des décisions
          consignées.
    \item Préparer la salle, l'annuaire, les relèves et la fin de crise.
\end{itemize}
\end{boxobjectifs}

\bigskip

% ============================================================================
\section*{Question directrice : pourquoi la réunion du 6 mars à 10 h 30 n'a-t-elle rien décidé~?}
% ============================================================================

Huit personnes, un directeur général au téléphone, aucun ordre du jour,
aucune personne chargée d'écrire. Chacun apporte ses informations, souvent
contradictoires~; la discussion porte sur les causes plutôt que sur les
décisions. À 11~h~10, on coupe enfin les liaisons. Le 7~mars, Mme~NKENG
ouvre une main courante, et la cellule commence à fonctionner. Ce chapitre,
sur deux séances, montre comment organiser ce qui a manqué le premier jour.

\medskip
\textbf{Plan du chapitre.}
\begin{enumerate}
    \item Deux niveaux (section~\ref{sec:c22-niveaux}).
    \item Les rôles (section~\ref{sec:c22-roles}).
    \item La main courante (section~\ref{sec:c22-main}).
    \item Le rythme de la crise (section~\ref{sec:c22-rythme}).
    \item La première heure (section~\ref{sec:c22-heure}).
    \item Salle, annuaire, relèves, fin (section~\ref{sec:c22-moyens}).
\end{enumerate}

\separateur

% ============================================================================
\section{Deux Niveaux}
\label{sec:c22-niveaux}
% ============================================================================

\begin{boxdef}
\termecle{Cellule de crise}~: équipe constituée à l'avance, réunie sur
décision de la direction, qui centralise l'information, prend ou prépare
les décisions, coordonne l'action et la communication pendant la crise.
\end{boxdef}

\begin{tableaucours}{L{3.0cm}L{4.6cm}L{4.6cm}}{%
    \tableauHeader{Niveau} &
    \tableauHeader{Rôle} &
    \tableauHeader{Chez LiZbiZ}
}
Stratégique (décisionnel) & Fixe les priorités, décide des arbitrages
(rançon, communication publique, clients), parle au nom de l'entreprise &
DG, DAF, directeur des opérations, juriste, communication \\
Opérationnel & Exécute, qualifie la situation, propose des options,
remonte l'information & Responsable informatique, prestataire, chefs
d'agence \\
\end{tableaucours}
\vspace{-0.8em}
\begin{center}{\small\itshape Tableau~22.1. Les deux niveaux de la cellule de crise}\end{center}

L'ISO~22361 décrit cette articulation entre le niveau qui décide et celui
qui agit~\cite{iso22361_2022}~; l'ISO~22320, pour les situations
d'urgence, insiste sur une chaîne de commandement claire et une
information partagée~\cite{iso22320_2018}.

\begin{boxalert}
\textbf{Violation typique.} Faire travailler les mêmes personnes aux deux
niveaux.

\smallskip
\textbf{Règle~:} M.~KOUNOU ne peut pas à la fois restaurer les serveurs et
siéger trois heures en cellule stratégique. Il y délègue un représentant
qui remonte la situation technique~; lui reste à la manœuvre.
\end{boxalert}

% ============================================================================
\section{Les Rôles}
\label{sec:c22-roles}
% ============================================================================

\begin{tableaucours}{L{3.6cm}L{5.2cm}L{3.4cm}}{%
    \tableauHeader{Rôle} &
    \tableauHeader{Mission} &
    \tableauHeader{Titulaire et suppléant}
}
Directeur de crise & Décide, arbitre, fixe le rythme & M.~ETEME~;
Mme~MBALLA \\
Coordinateur & Prépare les réunions, suit les actions, relance & Mme~MOUAFO~;
M.~NGUELE \\
Main courante & Consigne heure, fait, source, décision & Mme~NKENG~; une
assistante formée \\
Communication & Messages internes et externes, médias & Mme~ABOMO~; un
prestataire désigné \\
Juridique & Preuves, plainte, rançon, données personnelles & Me~TSALA~;
un second cabinet \\
Experts & Situation technique, options, délais & Représentant de
l'informatique, prestataire \\
Logistique & Salle, repas, transport, relèves & Services généraux \\
\end{tableaucours}
\vspace{-0.8em}
\begin{center}{\small\itshape Tableau~22.2. Les rôles de la cellule de crise de LiZbiZ}\end{center}

Chaque rôle a un \emph{suppléant} désigné et formé~: une crise dure
souvent plus longtemps que la résistance d'une personne, et le titulaire
peut être absent (en mars, le DG était sur la route de Yaoundé).

% ============================================================================
\section{La Main Courante}
\label{sec:c22-main}
% ============================================================================

\begin{boxdef}
\termecle{Main courante}~: registre chronologique, tenu en continu
pendant la crise, de tous les faits, informations, décisions et actions,
avec leur heure, leur source et leur auteur. Elle sert à piloter, à
transmettre lors des relèves, à rendre compte, et à apprendre.
\end{boxdef}

\begin{tableaucours}{L{1.4cm}L{1.6cm}L{4.6cm}L{2.4cm}L{2.0cm}}{%
    \tableauHeader{Heure} &
    \tableauHeader{Type} &
    \tableauHeader{Contenu} &
    \tableauHeader{Source ou auteur} &
    \tableauHeader{Échéance}
}
08~h~05 & Fait & Yaoundé et Bafoussam opérationnels, sans trace de
chiffrement & M.~KOUNOU & -- \\
08~h~20 & Décision & Ne pas payer la rançon tant que l'avis juridique n'est
pas rendu & M.~ETEME & 8~mars \\
08~h~25 & Action & Déposer plainte à la police judiciaire & Me~TSALA & 7~mars
12~h \\
08~h~40 & Information & Un client majeur demande un point écrit avant
midi & Mme~ABOMO & 11~h~30 \\
09~h~00 & Action & Diffuser le communiqué aux clients & Mme~ABOMO & 9~h~00 \\
\end{tableaucours}
\vspace{-0.8em}
\begin{center}{\small\itshape Tableau~22.3. Extrait de la main courante du 7 mars (Opération GAROUA)}\end{center}

\begin{boxprocedure}
\textbf{Tenir la main courante}
\begin{enumerate}
    \item Une ligne par entrée~: heure, type (fait, information, décision,
          action), contenu, source ou auteur, échéance.
    \item Distinguer le fait constaté de l'information rapportée~; noter la
          source.
    \item Relire chaque décision à voix haute avant de l'inscrire.
    \item Numéroter les actions et les suivre jusqu'à leur clôture.
    \item Ne rien effacer~; corriger par une nouvelle ligne datée.
    \item Conserver la main courante comme pièce du dossier.
\end{enumerate}
\end{boxprocedure}

% ============================================================================
\section{Le Rythme de la Crise}
\label{sec:c22-rythme}
% ============================================================================

La cellule se réunit à heures fixes (par exemple toutes les trois heures en
phase aiguë, deux fois par jour ensuite). Entre deux réunions, chacun agit.
Chaque réunion suit le même ordre.

\begin{boxprocedure}
\textbf{Ordre du jour d'un point de situation}
\begin{enumerate}
    \item Point de situation en trois colonnes~: ce que l'on sait, ce que
          l'on ne sait pas, ce que l'on fait (Chapitre~\ref{chap:crise}).
    \item Revue des actions de la réunion précédente.
    \item Décisions à prendre, avec les options et leurs conséquences.
    \item Communication~: messages à valider, destinataires.
    \item Heure du prochain point.
\end{enumerate}
\end{boxprocedure}

% ============================================================================
\section{La Première Heure}
\label{sec:c22-heure}
% ============================================================================

\begin{boxprocedure}
\textbf{Fiche réflexe du directeur de crise : la première heure}
\begin{enumerate}
    \item Réunir la cellule (physiquement ou à distance) et désigner la
          personne qui tient la main courante.
    \item Faire établir le premier point de situation.
    \item Décider des mesures conservatoires~: confiner (couper les
          liaisons), protéger les sauvegardes, préserver les preuves.
    \item Activer les plans~: continuité, reprise informatique.
    \item Désigner le porte-parole~; valider un premier message d'attente
          interne et externe.
    \item Prévenir le conseil juridique, l'assureur, le prestataire.
    \item Fixer l'heure du prochain point.
\end{enumerate}
\end{boxprocedure}

\begin{boxscenario}{Opération GAROUA. Le 6 mars avec une fiche réflexe}
\textbf{Constatation.} Le 6~mars, la cellule s'est réunie à 10~h~30, trois
heures après la découverte, sans main courante~; la coupure des liaisons a
été décidée à 11~h~10, le premier message externe le 7~mars à 9~h.

\medskip
\textbf{Analyse} (confiance moyenne). Avec la fiche, l'appel de 7~h~45 au
DG aurait suffi à déclencher la cellule à distance~: les liaisons
pouvaient être coupées vers 8~h, et un message d'attente publié dans la
matinée, avant la rumeur du soir.
\end{boxscenario}

% ============================================================================
\section{Salle, Annuaire, Relèves, Fin}
\label{sec:c22-moyens}
% ============================================================================

\begin{tableaucours}{L{3.0cm}L{9.2cm}}{%
    \tableauHeader{Moyen} &
    \tableauHeader{Exigence}
}
Salle & Une salle principale et une de repli, accessibles sans le réseau de
l'entreprise~; tableau, téléphones, alimentation de secours \\
Annuaire de crise & Imprimé, à jour~: cellule, suppléants, prestataires,
clients majeurs, assureur, autorités, médias \\
Moyens indépendants & Téléphones et messagerie hors du système compromis \\
Relèves & Passation écrite à partir de la main courante, toutes les 12~h
au plus \\
Fin de crise & Décidée et annoncée par le directeur de crise, selon des
critères fixés~; passage à une organisation allégée \\
\end{tableaucours}
\vspace{-0.8em}
\begin{center}{\small\itshape Tableau~22.4. Les moyens de la cellule de crise}\end{center}

\begin{boxobjet}{R-20 · Indisponibilité de la cellule de crise}
\textbf{Scénario~:} du fait de la messagerie hébergée sur le système
compromis, la cellule ne peut ni se joindre ni joindre les tiers.
\textbf{Cotation~:} V=2, I=3. \textbf{Traitement~:} annuaire imprimé,
groupe de messagerie hors du système, salle de repli.
\textbf{Propriétaire~:} Mme~MOUAFO. \textbf{Indicateur~:} date de la dernière
vérification de l'annuaire de crise.
\end{boxobjet}

\begin{boxmetier}{Directeur de crise}
\textbf{Décider de la fin.} Le directeur de crise ne laisse pas la crise
s'éteindre d'elle-même~: il annonce la fin, la passation à l'organisation
normale, et la date du retour d'expérience.
\end{boxmetier}

\separateur

\begin{boxresume}
\textbf{Ce qu'il faut retenir de ce chapitre~:}
\begin{itemize}
    \item \textbf{Deux niveaux}~: stratégique qui décide, opérationnel qui
          agit~; pas les mêmes personnes.
    \item \textbf{Des rôles nommés}, chacun avec un suppléant formé.
    \item \textbf{La main courante}~: heure, type, contenu, source,
          échéance~; on n'efface rien.
    \item \textbf{Un rythme fixe}, un ordre du jour constant.
    \item \textbf{La première heure} se joue avec une fiche réflexe~;
          salle, annuaire et relèves se préparent à froid.
\end{itemize}
\end{boxresume}

\begin{boxinfo}
\textbf{Points de contrôle de l'élève.} Je sais~: distinguer les deux
niveaux~; attribuer les rôles~; tenir une main courante~; conduire un
point de situation~; dérouler la fiche réflexe de la première heure.
\end{boxinfo}

\bigskip

\begin{boxexo}[title={\ding{45}\quad Exercice~22.1. Quel niveau~? \niveaubase}]
Rangez au niveau stratégique ou opérationnel~: décider de ne pas payer la
rançon~; restaurer la base de facturation~; valider le communiqué.
\end{boxexo}

\begin{boxexo}[title={\ding{45}\quad Exercice~22.2. Le type \niveaubase}]
Donnez le type (fait, information, décision, action) de chaque ligne du
Tableau~22.3.
\end{boxexo}

\begin{boxexo}[title={\ding{45}\quad Exercice~22.3. Les suppléants \niveaubase}]
Pourquoi chaque rôle a-t-il un suppléant~? Donnez l'exemple du 6~mars.
\end{boxexo}

\begin{boxexo}[title={\ding{45}\quad Exercice~22.4. Une main courante \niveaumoyen}]
Rédigez la main courante de la matinée du 6~mars, de 7~h~30 à 11~h~10, à
partir du Tableau~0.2.
\end{boxexo}

\begin{boxexo}[title={\ding{45}\quad Exercice~22.5. Corriger \niveaumoyen}]
À 10~h, on apprend que l'information de 8~h~05 était fausse~: Bafoussam
est touché. Comment le consigner dans la main courante~?
\end{boxexo}

\begin{boxexo}[title={\ding{45}\quad Exercice~22.6. Le rythme \niveaumoyen}]
Proposez le calendrier des points de situation du 6 au 9~mars.
\end{boxexo}

\begin{boxexo}[title={\ding{45}\quad Exercice~22.7. La relève \niveaumoyen}]
Rédigez la fiche de relève remise par Mme~NKENG à sa suppléante le
7~mars à 20~h.
\end{boxexo}

\begin{boxexo}[title={\ding{45}\quad Exercice~22.8. L'annuaire \niveaumoyen}]
Dressez la liste des entrées de l'annuaire de crise de LiZbiZ, par
catégorie.
\end{boxexo}

\begin{boxexo}[title={\ding{45}\quad Exercice~22.9. Le plan de gestion de crise \niveauexpert}]
Rédigez le plan de gestion de crise de LiZbiZ~: déclenchement, cellule,
rôles, salle, moyens, rythme, fin.
\end{boxexo}

\begin{boxexo}[title={\ding{45}\quad Exercice~22.10. La cellule bancaire \niveauexpert}]
Composez la cellule de crise de la BANQUE SANAGA pour une indisponibilité du
système central, en tenant compte des autorités de supervision et des
systèmes de paiement.
\end{boxexo}

\bigskip

\begin{boxinfo}
\textbf{Pour aller plus loin.}
\begin{itemize}
    \item ISO~22361~\cite{iso22361_2022}~; ISO~22320~\cite{iso22320_2018}.
\end{itemize}
\end{boxinfo}

\begin{boxinfo}
\textbf{Transition.} La cellule organisée, il reste à bien décider. Le
Chapitre~\ref{chap:decider} traite de la décision en situation de crise.
\end{boxinfo}
